% problem-000464 5 MAX相层状晶体结构
\section*{5. MAX相层状晶体结构}
MAX（M代表过渡⾦属元素，A代表主族元素，X代表碳或氮）相是⼀类备受关注的新型陶瓷材料。由于独特的层状晶体结构，其具有⾃润滑、⾼韧性、可导电等性能，可作为⾼温结构材料、电极材料和化学防腐材料。某MAX相材料含有Ti，Al，N3种原⼦，属六⽅晶系，Ti原⼦的堆积⽅式为…BACBBCABBACBBCAB…，其中A、B、C 都是密置单层。N 原⼦占据所有的正⼋⾯体空隙，⽽Al 原⼦占据⼀半的三棱柱空隙。如果Ti 原⼦层上下同时接触N和Al原⼦，则沿着晶胞�轴⽅向，Al和N原⼦的投影重合。

\subsection*{5-1}
写出该化合物的化学式，及每个正当晶胞中的原⼦种类和个数。

\subsection*{5-2}
沿着晶胞�轴⽅向，画⼀条同时含有Al和N原⼦的直线，标出直线上的原⼦排列（⽆需考虑原⼦间距离，直线上总原⼦数不少于10个。Al，Ti，N分别⽤○，△，□表⽰）。

\subsection*{5-3}
已知 Ti，N 原⼦之间的平均键⻓为 210.0 pm，Ti，Al 原⼦之间的平均键⻓为 281.8 pm，估算晶体的理论密度（原⼦量：Ti: 47.87，Al: 26.98，N: 14.01， $N _ { \mathrm { A } } = 6 . 0 2 { \times } 1 0 ^ { 2 3 } \mathrm { m o l ^ { - 1 } } )$ o

\subsection*{5-4}
晶粒尺⼨会影响上述材料的性质，所以⾼温制备时⼀般通过延⻓保温时间来增加晶粒尺⼨ $\mathbf { \rho } _ { \circ }$ 。判断常温下晶粒⽣⻓过程的熵变、焓变和⾃由能变化的正负，并从化学热⼒学⻆度判断常温下该晶粒⽣⻓过程是否⾃发。

\subsection*{5-5}
以上描述均针对完美晶体。⼀般情况下，晶粒中会出现缺陷。从热⼒学⻆度证明：对于⾜够⼤的晶体，出现缺陷是⾃发的。
